%% ma: L10-B01-0033
%% lop: 10
%% chuong: 1
%% bai: 1
%% dang: 10.1.6
%% loai: TLN
%% muc: TH
%% dap_an: "4"
%% nguon: "Ảnh giáo viên gửi 10/10/2026 (ThS. Nguyễn Hoàng Việt, Chương 1 Mệnh đề), B-Đề 2, Phần 3 Câu 17"
%% kien_thuc: "Bài 1 (mệnh đề kéo theo, mệnh đề tương đương)"
%% trang_thai: chinh-thuc
%% ghi_chu: "Giáo viên đã duyệt 10/10/2026. Mệnh đề thứ hai và thứ tư của đề gốc giống hệt nhau"
\begin{ex}
Cho các mệnh đề sau. Hãy cho biết có bao nhiêu mệnh đề đúng?
\begin{itemize}
\item $\forall n\in\mathbb N,\ n^2$ chia hết cho $7\Rightarrow n$ chia hết cho $7$.
\item $\forall n\in\mathbb N,\ n^2$ chia hết cho $5\Rightarrow n$ chia hết cho $5$.
\item Nếu tam giác $ABC$ không phải là tam giác đều thì tam giác đó có ít nhất một góc nhỏ hơn $60^\circ$.
\item $\forall n\in\mathbb N,\ n^2$ chia hết cho $5\Rightarrow n$ chia hết cho $5$.
\end{itemize}
\shortans{4}
\loigiai{$7$ và $5$ là các số nguyên tố nên nếu $p\mid n^2$ thì $p\mid n$: mệnh đề thứ nhất, thứ hai và thứ tư đúng. Mệnh đề thứ ba đúng: ba góc có tổng $180^\circ$, nếu không có góc nào nhỏ hơn $60^\circ$ thì cả ba góc đều $\ge60^\circ$ nên cùng bằng $60^\circ$, tam giác đều. Vậy có $4$ mệnh đề đúng.}
\end{ex}
